\chapter{Fibrations de Lefschetz}

\section{Introduction aux fibrations de Lefschetz}

Les fibrations de Lefschetz constituent un pont essentiel entre la géométrie algébrique complexe et la topologie différentielle des variétés de dimension 4. Initialement introduites par Lefschetz dans le cadre de l'étude des variétés algébriques, elles ont connu un regain d'intérêt considérable depuis les travaux de Donaldson \cite{Donaldson99} montrant que toute variété symplectique de dimension 4 admet une telle fibration.

\subsection{Définition locale}

Commençons par décrire le modèle local d'un point critique de type Lefschetz.

\begin{definition}
Soit $X$ une variété différentiable de dimension 4 et $\Sigma$ une surface de Riemann. Une application différentiable $\pi: X \to \Sigma$ admet un \emph{point critique de type Lefschetz} en $x \in X$ s'il existe des coordonnées locales complexes $(z_1, z_2)$ au voisinage de $x$ et une coordonnée locale $t$ au voisinage de $\pi(x)$ telles que
\[
\pi(z_1, z_2) = z_1^2 + z_2^2.\]
\end{definition}

Ce modèle local est fondamental. La fibre $\pi^{-1}(t)$ pour $t \neq 0$ est une surface de Riemann lisse ; lorsque $t = 0$, la fibre présente un point double ordinaire (un nœud). La transformation monodromique associée au tour autour de la valeur critique est un \emph{twist de Dehn} le long du cycle évanouissant.

\subsection{Définition globale}

\begin{definition}
Soit $X$ une variété différentiable compacte connexe de dimension 4 et $\Sigma$ une surface de Riemann compacte connexe. Une \emph{fibration de Lefschetz} est une application $\pi: X \to \Sigma$ vérifiant les propriétés suivantes :
\begin{enumerate}
\item $\pi$ est propre et à fibres connexes ;
\item L'ensemble $\Delta \subset X$ des points critiques de $\pi$ est fini, et l'ensemble $B = \pi(\Delta) \subset \Sigma$ est l'ensemble des valeurs critiques ;
\item En chaque point critique $x \in \Delta$, $\pi$ admet un modèle local de la forme $z_1^2 + z_2^2$ comme ci-dessus ;
\item $\pi$ est une submersion en dehors de $\Delta$.
\end{enumerate}
La fibre régulière $\pi^{-1}(b)$ pour $b \in \Sigma \setminus B$ est une surface de Riemann orientée, dont on note $g$ le genre. On parle alors de \emph{fibration de Lefschetz de genre $g$}.
\end{definition}

\begin{remarque}
La condition de connexité des fibres assure que la fibration est relativement minimale au sens topologique : on ne peut pas contracter de composantes de fibres sans créer de singularités.
\end{remarque}

\subsection{Premières propriétés}

Soit $\pi: X \to \Sigma$ une fibration de Lefschetz de genre $g$. Notons $k = |B|$ le nombre de valeurs critiques. Alors $X$ admet une décomposition cellulaire naturelle : on peut reconstruire $X$ à partir de la fibre régulière $F$ en attachant $k$ anses de dimension 2 le long des cycles évanouissants. Cette construction est l'analogue topologique de la description des surfaces elliptiques par leurs fibres singulières.

La caractéristique d'Euler de $X$ se calcule par la formule
\[
\chi(X) = \chi(F) \cdot \chi(\Sigma) + k,
\]
où $\chi(F) = 2-2g$ et $\chi(\Sigma) = 2-2g(\Sigma)$. Cette formule reflète le fait que chaque point critique contribue un supplément de 1 à la caractéristique d'Euler (car le modèle local $z_1^2+z_2^2$ a une singularité isolée d'indice 1).

\section{Groupe de mapping class et twists de Dehn}

\subsection{Le groupe de mapping class}

\begin{definition}
Soit $\Sigma_g$ une surface de Riemann compacte connexe orientée de genre $g$. Le \emph{groupe de mapping class} $\Gamma_g$ est le groupe
\[
\Gamma_g = \mathrm{Diff}^+(\Sigma_g) / \mathrm{Diff}_0(\Sigma_g),
\]
où $\mathrm{Diff}^+(\Sigma_g)$ est le groupe des difféomorphismes préservant l'orientation de $\Sigma_g$ et $\mathrm{Diff}_0(\Sigma_g)$ est le sous-groupe des difféomorphismes isotopes à l'identité.
\end{definition}

Le groupe $\Gamma_g$ est un objet discret qui encode la structure globale de la surface. Pour $g=0$, $\Gamma_0$ est trivial (tout difféomorphisme de la sphère est isotope à une rotation, et on peut déformer une rotation en l'identité). Pour $g=1$, on a un isomorphisme classique :
\[
\Gamma_1 \cong SL(2,\Z).
\]
Pour $g \ge 2$, $\Gamma_g$ est un groupe infini, non abélien, avec une présentation explicite due à Dehn.

\subsection{Twists de Dehn}

Les twists de Dehn sont les générateurs fondamentaux du groupe de mapping class.

\begin{definition}
Soit $\gamma \subset \Sigma_g$ une courbe simple fermée (non contractile). Un \emph{twist de Dehn} le long de $\gamma$ est un difféomorphisme $t_\gamma: \Sigma_g \to \Sigma_g$ défini de la manière suivante : dans un voisinage annulaire $A \cong S^1 \times [-1,1]$ de $\gamma$ paramétré par $(\theta, r)$, on pose
\[
t_\gamma(\theta, r) = (\theta + 2\pi r, r),
\]
et on prolonge par l'identité en dehors de $A$.
\end{definition}

Intuitivement, on coupe la surface le long de $\gamma$, on fait tourner l'un des bords d'un tour complet, puis on recolle. La classe d'isotopie de $t_\gamma$ ne dépend que de la classe d'isotopie de $\gamma$.

\begin{theoreme}[Dehn]
Le groupe $\Gamma_g$ est engendré par les twists de Dehn le long d'un nombre fini de courbes simples.
\end{theoreme}

Dans le cas du genre 1, un twist de Dehn le long d'un méridien (ou d'une longitude) correspond à la matrice
\[
\begin{pmatrix} 1 & 1 \\ 0 & 1 \end{pmatrix} \in SL(2,\Z)
\]
pour un choix convenable de base de l'homologie. Les deux twists standards $U$ et $V$ (le long d'un méridien et d'une longitude) satisfont la relation
\[
(UV)^6 = I
\]
dans $SL(2,\Z)$, ce qui correspond à la présentation classique du groupe modulaire.

\section{Monodromie d'une fibration de Lefschetz}

\subsection{Définition de la monodromie}

Soit $\pi: X \to \Sigma$ une fibration de Lefschetz de genre $g$, avec valeurs critiques $B = \{p_1, \dots, p_k\} \subset \Sigma$. Choisissons un point base $b_0 \in \Sigma \setminus B$ et notons $F_0 = \pi^{-1}(b_0)$ la fibre régulière.

Pour tout lacet $\gamma: [0,1] \to \Sigma \setminus B$ basé en $b_0$, le transport parallèle le long de $\gamma$ (défini par relèvement des chemins) induit un difféomorphisme $h_\gamma: F_0 \to F_0$. Ce difféomorphisme est défini à isotopie près (car il dépend du choix des relèvements).

\begin{definition}
La \emph{représentation de monodromie} de la fibration de Lefschetz est le morphisme
\[
\rho: \pi_1(\Sigma \setminus B, b_0) \longrightarrow \Gamma_g
\]
défini par $\rho([\gamma]) = [h_\gamma]$, où $[\cdot]$ désigne la classe d'isotopie.
\end{definition}

\begin{theoreme}
La représentation de monodromie est un invariant complet de la fibration de Lefschetz à difféomorphisme près, préservant la fibration.
\end{theoreme}

\subsection{Monodromie locale autour d'un point critique}

Soit $p \in B$ une valeur critique. Considérons un petit lacet $\gamma$ qui tourne une fois autour de $p$ (dans le sens positif) et n'enlace aucun autre point critique. La monodromie locale $\rho(\gamma)$ est un élément particulier de $\Gamma_g$.

\begin{theoreme}
Soit $x$ l'unique point critique au-dessus de $p$. Le cycle évanouissant $\delta \subset F_0$ est la courbe simple qui se contracte sur $x$ lorsque l'on s'approche de $p$. Alors $\rho(\gamma)$ est le twist de Dehn $t_\delta$ le long de $\delta$.
\end{theoreme}

Ce résultat fondamental relie la géométrie locale des points critiques à la structure globale du groupe de mapping class. La courbe $\delta$ est appelée \emph{cycle évanouissant} car elle se contracte sur le point singulier lorsque la fibre dégénère.

\subsection{Relations de monodromie globale}

Supposons maintenant que $\Sigma = \CP^1$ (le cas le plus fréquent dans les applications). Alors $\pi_1(\CP^1 \setminus B, b_0)$ est un groupe libre à $k$ générateurs $\gamma_1, \dots, \gamma_k$, où $\gamma_i$ est un lacet qui tourne autour de $p_i$ (et contourne les autres points critiques).

La condition de compatibilité globale provient du fait qu'un lacet entourant tous les points critiques est contractile dans $\CP^1$ (car $\CP^1$ est simplement connexe). On obtient donc la relation
\[
\rho(\gamma_1) \cdots \rho(\gamma_k) = I \quad \text{dans } \Gamma_g.
\]

Ainsi, une fibration de Lefschetz sur $\CP^1$ est entièrement décrite par une factorisation de l'identité en produit de twists de Dehn :
\[
t_{\delta_1} \cdots t_{\delta_k} = I,
\]
où chaque $\delta_i$ est le cycle évanouissant associé au point critique $p_i$.

\section{Lien avec les surfaces elliptiques}

\subsection{D'une surface elliptique à une fibration de Lefschetz}

Soit $\pi: S \to C$ une surface elliptique relativement minimale. Comme surface complexe, $S$ est en particulier une variété différentiable de dimension 4, et $\pi$ est une application différentiable. En lissant éventuellement certaines singularités, on obtient :

\begin{proposition}
Toute surface elliptique relativement minimale $\pi: S \to C$ définit, après perturbation générique si nécessaire, une fibration de Lefschetz de genre 1.
\end{proposition}

La perturbation est nécessaire car dans une surface elliptique complexe, les points singuliers des fibres ne sont pas nécessairement des points critiques de type Lefschetz au sens différentiable (ils peuvent correspondre à des singularités plus compliquées, comme des cusps). Cependant, une petite déformation permet de remplacer chaque fibre singulière par une configuration de fibres de type $I_n$ qui sont de vraies singularités de Lefschetz.

\begin{remarque}
La classification de Kodaira trouve ici son pendant topologique : chaque type de fibre singulière se résout, après perturbation, en une collection de fibres $I_1$ (singularités de Lefschetz). Par exemple, une fibre $I_n$ donne $n$ points critiques distincts, tandis qu'une fibre $II$ donne un point critique de type cusp qui doit être déformé.
\end{remarque}

\subsection{Réciproque : quand une fibration de Lefschetz est-elle algébrique ?}

Toutes les fibrations de Lefschetz ne proviennent pas d'une structure complexe. La question de savoir quand une fibration de Lefschetz de genre 1 est \emph{algébrique} (c'est-à-dire provient d'une surface elliptique complexe) est délicate.

\begin{theoreme}
Une fibration de Lefschetz de genre 1 sur $\CP^1$ provient d'une surface elliptique complexe (après éventuelle perturbation) si et seulement si sa monodromie peut être factorisée en twists de Dehn correspondant à des cycles évanouissants qui sont des courbes simples non séparantes et si les paramètres complexes peuvent être choisis cohéremment (condition de périodes).
\end{theoreme}

En pratique, les fibrations de Lefschetz que nous considérerons dans ce mémoire seront toujours issues de surfaces elliptiques complexes, et nous nous concentrerons sur l'information topologique contenue dans la monodromie.

\section{Exemples fondamentaux}

\subsection{Le cas du genre 1 et $SL(2,\Z)$}

Pour $g=1$, la fibre régulière est un tore $\Sigma_1$. Le groupe de mapping class $\Gamma_1$ s'identifie à $SL(2,\Z)$ agissant sur l'homologie $H_1(\Sigma_1, \Z) \cong \Z^2$. Un twist de Dehn le long d'une courbe simple non contractile (un méridien ou une longitude) correspond à une matrice unipotente :
\[
U = \begin{pmatrix} 1 & 1 \\ 0 & 1 \end{pmatrix}, \quad V = \begin{pmatrix} 1 & 0 \\ -1 & 1 \end{pmatrix}.
\]

Ces deux matrices engendrent $SL(2,\Z)$ et satisfont la relation fondamentale
\[
(UV)^6 = I
\]
(dans $SL(2,\Z)$, ce produit vaut $-I$, mais dans $PSL(2,\Z)$ on a bien l'ordre 6).

\subsection{Fibration triviale et produit}

L'exemple le plus simple est le produit $X = \Sigma_1 \times \CP^1$ avec la projection sur le second facteur. Cette fibration n'a pas de points critiques ( $k=0$ ), et sa monodromie est triviale. C'est une surface elliptique (produit) mais elle n'est pas relativement minimale au sens algébrique car elle contient des courbes exceptionnelles.

\subsection{Fibration de Lefschetz à 12 singularités}

La surface $E(1)$ construite au chapitre précédent donne, après perturbation des fibres $I_1$, une fibration de Lefschetz de genre 1 avec $k=12$ points critiques. La monodromie est donnée par 12 twists de Dehn dont le produit est l'identité. Nous étudierons cet exemple en détail au Chapitre 4.

\subsection{Construction par surgeries}

Une méthode générale pour construire des fibrations de Lefschetz est la \emph{soudure de tores} (mapping torus) : si l'on se donne une factorisation de l'identité en twists de Dehn, on peut reconstruire la 4-variété correspondante en recollant des morceaux standards. Cette construction est à la base de l'approche combinatoire de Gompf et Stipsicz \cite{GS99}.

\section{Discussion : invariants topologiques}

Pour une fibration de Lefschetz $\pi: X \to \CP^1$ de genre $g$ avec $k$ points critiques, on a les formules suivantes :

\begin{itemize}
\item Caractéristique d'Euler : $\chi(X) = (2-2g) \cdot 2 + k = 4-4g + k$.
\item Signature : $\sigma(X)$ peut se calculer à partir de la monodromie via la formule de Meyer \cite{Meyer72} qui exprime la signature d'un fibré en tores en termes du cocycle de Meyer sur $SL(2,\Z)$.
\item Genre de la fibre régulière : $g$ est déterminé par la monodromie (car les cycles évanouissants sont des courbes dans $\Sigma_g$).
\end{itemize}

Dans le cas du genre 1, la signature est donnée par
\[
\sigma(X) = -\frac{1}{3} \sum_{i=1}^k \varepsilon_i,
\]
où $\varepsilon_i = \pm 1$ selon que le twist est positif ou négatif (correspondant à $U$ ou $U^{-1}$). Cette formule sera essentielle dans l'étude des exemples.
